\section{行列式}

\subsection{排列与逆序数}

\begin{mydef}{排列与逆序}{permutation and inversion}
    由 $1,2,\ldots,n$ 组成的一个有序数组称为一个 $n$ 级排列，记作 $p_1p_2\ldots p_n$。\\
    在一个排列中，如果一对数的前后位置与大小顺序相反，即前面的数大于后面的数，
    则称它们构成一个\textbf{逆序}（inversion）。\\
    一个排列中逆序的总数称为该排列的\textbf{逆序数}，记作 $\tau(p_1p_2\ldots p_n)$。
\end{mydef}

\begin{mydef}{奇排列与偶排列}{even and odd permutation}
    逆序数为偶数的排列称为\textbf{偶排列}，逆序数为奇数的排列称为\textbf{奇排列}。
\end{mydef}

\begin{conclusion}{对换改变排列的奇偶性}{transposition changes parity}
    将一个排列中的两个元素对换位置（称为一次\textbf{对换}），排列的奇偶性发生改变。\\
    由此可知，在全部 $n!$ 个 $n$ 级排列中，奇排列与偶排列各占一半，各有 $\frac{n!}{2}$ 个。
\end{conclusion}

\begin{formula}{逆序数的计算}{computing inversion number}
    设 $n$ 级排列 $p_1p_2\ldots p_n$，则
    \[
    \tau(p_1p_2\ldots p_n) = \sum_{i=1}^{n-1} t_i
    \]
    其中 $t_i$ 表示在 $p_i$ 后面且比 $p_i$ 小的数的个数。\\
    \boxed{\textbf{例}} 排列 $53241$ 的逆序数：$5$ 后面有 $3,2,4,1$ 四个比它小，$3$ 后面有 $2,1$ 两个，
    $2$ 后面有 $1$ 一个，$4$ 后面有 $1$ 一个，故 $\tau(53241)=4+2+1+1=8$（偶排列）。
\end{formula}

\subsection{行列式的定义}

\begin{mydef}{$n$ 阶行列式的排列定义}{determinant via permutations}
    由 $n^2$ 个数排成 $n$ 行 $n$ 列，记作
    \[
    D = \begin{vmatrix}
    a_{11} & a_{12} & \cdots & a_{1n} \\
    a_{21} & a_{22} & \cdots & a_{2n} \\
    \vdots & \vdots & \ddots & \vdots \\
    a_{n1} & a_{n2} & \cdots & a_{nn}
    \end{vmatrix}
    \]
    称为 $n$ 阶行列式，它等于所有取自不同行不同列的 $n$ 个元素的乘积的代数和：
    \[
    \det(A) = \sum_{(p_1p_2\ldots p_n)} (-1)^{\tau(p_1p_2\ldots p_n)} a_{1p_1}a_{2p_2}\cdots a_{np_n}
    \]
    其中求和取遍所有 $n$ 级排列 $p_1p_2\ldots p_n$，共有 $n!$ 项。
\end{mydef}

\begin{formula}{二阶与三阶行列式}{2nd and 3rd order determinants}
    \textbf{二阶行列式}：
    \[
    \begin{vmatrix}
    a_{11} & a_{12} \\
    a_{21} & a_{22}
    \end{vmatrix}
    = a_{11}a_{22} - a_{12}a_{21}
    \]
    \textbf{三阶行列式}（对角线法则）：
    \[
    \begin{vmatrix}
    a_{11} & a_{12} & a_{13} \\
    a_{21} & a_{22} & a_{23} \\
    a_{31} & a_{32} & a_{33}
    \end{vmatrix}
    = a_{11}a_{22}a_{33} + a_{12}a_{23}a_{31} + a_{13}a_{21}a_{32}
    - a_{13}a_{22}a_{31} - a_{12}a_{21}a_{33} - a_{11}a_{23}a_{32}
    \]
    \boxed{\text{注}} 对角线法则仅适用于二阶和三阶行列式，对 $n \geq 4$ 不成立。
\end{formula}

\subsection{行列式的基本性质}

\begin{conclusion}{转置不变性}{transpose invariance}
    行列式与它的转置行列式相等：
    \[
    \det(A^{T}) = \det(A)
    \]
    即行列式中行与列的地位是对称的，凡对行成立的性质对列也成立。
\end{conclusion}

\begin{conclusion}{交换行（列）变号}{row swap flips sign}
    交换行列式的两行（或两列），行列式变号：
    \[
    \begin{vmatrix}
    \cdots & a_{i1} & \cdots \\
    \cdots & a_{j1} & \cdots
    \end{vmatrix}
    = -\,
    \begin{vmatrix}
    \cdots & a_{j1} & \cdots \\
    \cdots & a_{i1} & \cdots
    \end{vmatrix}
    \]
    \textbf{推论}：若行列式有两行（列）完全相同，则行列式为零。
\end{conclusion}

\begin{conclusion}{公因子提取}{common factor extraction}
    行列式中某一行（列）的所有元素的公因子可以提到行列式符号外面：
    \[
    \begin{vmatrix}
    a_{11} & \cdots & k a_{1j} & \cdots & a_{1n} \\
    a_{21} & \cdots & k a_{2j} & \cdots & a_{2n} \\
    \vdots & \vdots & \vdots & \ddots & \vdots \\
    a_{n1} & \cdots & k a_{nj} & \cdots & a_{nn}
    \end{vmatrix}
    = k\,
    \begin{vmatrix}
    a_{11} & \cdots & a_{1j} & \cdots & a_{1n} \\
    a_{21} & \cdots & a_{2j} & \cdots & a_{2n} \\
    \vdots & \vdots & \vdots & \ddots & \vdots \\
    a_{n1} & \cdots & a_{nj} & \cdots & a_{nn}
    \end{vmatrix}
    \]
    \textbf{推论}：若行列式有两行（列）成比例，则行列式为零。\\
    特别地，若行列式中某一行（列）全为零，则行列式为零。
\end{conclusion}

\begin{conclusion}{行列式的线性性质}{linearity of determinant}
    若行列式的某一行（列）的元素是两组数之和，则此行列式等于两个行列式之和，
    这两个行列式分别以这两组数为该行（列）的元素，其余行（列）不变：
    \[
    \begin{vmatrix}
    a_{11} & \cdots & b_{1j} + c_{1j} & \cdots & a_{1n} \\
    \vdots & \vdots & \vdots & \ddots & \vdots \\
    a_{n1} & \cdots & b_{nj} + c_{nj} & \cdots & a_{nn}
    \end{vmatrix}
    =
    \begin{vmatrix}
    a_{11} & \cdots & b_{1j} & \cdots & a_{1n} \\
    \vdots & \vdots & \vdots & \ddots & \vdots \\
    a_{n1} & \cdots & b_{nj} & \cdots & a_{nn}
    \end{vmatrix}
    +
    \begin{vmatrix}
    a_{11} & \cdots & c_{1j} & \cdots & a_{1n} \\
    \vdots & \vdots & \vdots & \ddots & \vdots \\
    a_{n1} & \cdots & c_{nj} & \cdots & a_{nn}
    \end{vmatrix}
    \]
\end{conclusion}

\begin{conclusion}{行列式的倍加变换}{row addition leaves determinant unchanged}
    把行列式的某一行（列）的倍数加到另一行（列）上，行列式不变：
    \[
    \begin{vmatrix}
    \vdots & \vdots & \vdots \\
    a_{i1} & \cdots & a_{in} \\
    \vdots & \vdots & \vdots \\
    a_{j1} + k a_{i1} & \cdots & a_{jn} + k a_{in} \\
    \vdots & \vdots & \vdots
    \end{vmatrix}
    =
    \begin{vmatrix}
    \vdots & \vdots & \vdots \\
    a_{i1} & \cdots & a_{in} \\
    \vdots & \vdots & \vdots \\
    a_{j1} & \cdots & a_{jn} \\
    \vdots & \vdots & \vdots
    \end{vmatrix}
    \]
    这是行列式计算中最常用的性质，用于将行列式化为上三角形式。
\end{conclusion}

\boxed{\textbf{例}} 利用性质计算行列式：
\[
\begin{vmatrix}
2 & 1 & 3 \\
4 & 3 & 5 \\
6 & 5 & 5
\end{vmatrix}
\xrightarrow{r_2 - 2r_1}
\begin{vmatrix}
2 & 1 & 3 \\
0 & 1 & -1 \\
6 & 5 & 5
\end{vmatrix}
\xrightarrow{r_3 - 3r_1}
\begin{vmatrix}
2 & 1 & 3 \\
0 & 1 & -1 \\
0 & 2 & -4
\end{vmatrix}
\xrightarrow{r_3 - 2r_2}
\begin{vmatrix}
2 & 1 & 3 \\
0 & 1 & -1 \\
0 & 0 & -2
\end{vmatrix}
= 2 \cdot 1 \cdot (-2) = -4
\]

\subsection{行列式按行（列）展开}

\begin{mydef}{余子式与代数余子式}{minor and cofactor}
    在 $n$ 阶行列式中，划去元素 $a_{ij}$ 所在的第 $i$ 行和第 $j$ 列后，
    剩下的元素按原来的次序组成的 $n-1$ 阶行列式称为元素 $a_{ij}$ 的\textbf{余子式}，记作 $M_{ij}$。\\
    称 $A_{ij} = (-1)^{i+j} M_{ij}$ 为元素 $a_{ij}$ 的\textbf{代数余子式}（cofactor）。
\end{mydef}

\begin{conclusion}{行列式按行（列）展开定理}{cofactor expansion theorem}
    行列式等于它的任一行（列）的各元素与其对应的代数余子式乘积之和：
    \[
    \det(A) = \sum_{j=1}^{n} a_{ij} A_{ij} \quad (\text{按第 } i \text{ 行展开})
    \]
    \[
    \det(A) = \sum_{i=1}^{n} a_{ij} A_{ij} \quad (\text{按第 } j \text{ 列展开})
    \]
    此外，一行（列）的元素与另一行（列）的代数余子式的乘积之和为零：
    \[
    \sum_{k=1}^{n} a_{ik} A_{jk} = 0 \quad (i \neq j), \qquad
    \sum_{k=1}^{n} a_{ki} A_{kj} = 0 \quad (i \neq j)
    \]
\end{conclusion}

\boxed{\textbf{例}} 按第一行展开计算：
\[
\begin{aligned}
\begin{vmatrix}
2 & 1 & 0 \\
3 & 4 & 5 \\
0 & 6 & 7
\end{vmatrix}
&=2\begin{vmatrix}4&5\\6&7\end{vmatrix}
-\begin{vmatrix}3&5\\0&7\end{vmatrix}\\
&=2(28-30)-21=-25.
\end{aligned}
\]

\begin{conclusion}{范德蒙德行列式}{Vandermonde determinant}
    $n$ 阶范德蒙德行列式：
    \[
    V_n =
    \begin{vmatrix}
    1 & 1 & 1 & \cdots & 1 \\
    x_1 & x_2 & x_3 & \cdots & x_n \\
    x_1^2 & x_2^2 & x_3^2 & \cdots & x_n^2 \\
    \vdots & \vdots & \vdots & \ddots & \vdots \\
    x_1^{n-1} & x_2^{n-1} & x_3^{n-1} & \cdots & x_n^{n-1}
    \end{vmatrix}
    = \prod_{1 \leq i < j \leq n} (x_j - x_i)
    \]
    即 $V_n = (x_2-x_1)(x_3-x_1)\cdots(x_n-x_1)(x_3-x_2)\cdots(x_n-x_{n-1})$。\\
    \textbf{推论}：范德蒙德行列式为零的充要条件是 $x_1,x_2,\ldots,x_n$ 中至少有两个相等。
\end{conclusion}

\boxed{\textbf{例}} 计算四阶范德蒙德行列式：
\[
\begin{vmatrix}
1 & 1 & 1 & 1 \\
2 & 3 & -1 & 4 \\
4 & 9 & 1 & 16 \\
8 & 27 & -1 & 64
\end{vmatrix}
\begin{aligned}
&=(3-2)(-1-2)(4-2)\\
&\qquad\cdot(-1-3)(4-3)\bigl(4-(-1)\bigr)\\
&=1\cdot(-3)\cdot2\cdot(-4)\cdot1\cdot5=120
\end{aligned}
\]

\begin{conclusion}{三角行列式}{triangular determinant}
    \textbf{上三角行列式}（主对角线以下全为零）：
    \[
    \begin{vmatrix}
    a_{11} & a_{12} & \cdots & a_{1n} \\
    0 & a_{22} & \cdots & a_{2n} \\
    \vdots & \vdots & \ddots & \vdots \\
    0 & 0 & \cdots & a_{nn}
    \end{vmatrix}
    = a_{11}a_{22}\cdots a_{nn}
    \]
    \textbf{下三角行列式}（主对角线以上全为零）：
    \[
    \begin{vmatrix}
    a_{11} & 0 & \cdots & 0 \\
    a_{21} & a_{22} & \cdots & 0 \\
    \vdots & \vdots & \ddots & \vdots \\
    a_{n1} & a_{n2} & \cdots & a_{nn}
    \end{vmatrix}
    = a_{11}a_{22}\cdots a_{nn}
    \]
    对角行列式的值也等于主对角线元素的乘积。
\end{conclusion}

\begin{conclusion}{分块三角行列式}{block triangular determinant}
    设 $A$ 为 $m$ 阶方阵，$B$ 为 $n$ 阶方阵，则
    \[
    \begin{vmatrix}
    A & C \\
    0 & B
    \end{vmatrix}
    = \det(A) \cdot \det(B), \qquad
    \begin{vmatrix}
    A & 0 \\
    C & B
    \end{vmatrix}
    = \det(A) \cdot \det(B)
    \]
    特别地，对于分块对角矩阵：
    \[
    \begin{vmatrix}
    A & 0 \\
    0 & B
    \end{vmatrix}
    = \det(A) \cdot \det(B)
    \]
\end{conclusion}

\begin{conclusion}{行列式的乘法公式}{multiplicative property of determinants}
    设 $A$ 和 $B$ 都是 $n$ 阶方阵，则
    \[
    \det(AB) = \det(A) \cdot \det(B)
    \]
    由此可推出：$\det(A^k) = [\det(A)]^k$，$\det(kA) = k^n \det(A)$。
\end{conclusion}

\subsection{克莱姆法则}

\begin{conclusion}{克莱姆法则}{Cramer's rule}
    对于含有 $n$ 个方程 $n$ 个未知数的线性方程组
    \[
    \begin{cases}
    a_{11}x_1 + a_{12}x_2 + \cdots + a_{1n}x_n = b_1 \\
    a_{21}x_1 + a_{22}x_2 + \cdots + a_{2n}x_n = b_2 \\
    \qquad \qquad \qquad \cdots \cdots \\
    a_{n1}x_1 + a_{n2}x_2 + \cdots + a_{nn}x_n = b_n
    \end{cases}
    \]
    若系数行列式 $D = \det(a_{ij}) \neq 0$，则方程组有唯一解：
    \[
    x_1 = \frac{D_1}{D},\; x_2 = \frac{D_2}{D},\; \ldots,\; x_n = \frac{D_n}{D}
    \]
    其中 $D_j$ 是将 $D$ 中第 $j$ 列替换为常数项列 $(b_1,b_2,\ldots,b_n)^{T}$ 所得的行列式。
\end{conclusion}

\boxed{\textbf{例}} 用克莱姆法则解方程组：
\[
\begin{cases}
2x_1 + x_2 - x_3 = 1 \\
x_1 - x_2 + x_3 = 2 \\
3x_1 + 2x_2 - x_3 = 0
\end{cases}
\]
先计算系数行列式：
\[
D = \begin{vmatrix}
2 & 1 & -1 \\
1 & -1 & 1 \\
3 & 2 & -1
\end{vmatrix}
= 2(1-2) - 1((-1)-3) + (-1)(2+3) = -2 + 4 - 5 = -3 \neq 0
\]
\[
D_1 = \begin{vmatrix}
1 & 1 & -1 \\
2 & -1 & 1 \\
0 & 2 & -1
\end{vmatrix} = -3,\;
D_2 = \begin{vmatrix}
2 & 1 & -1 \\
1 & 2 & 1 \\
3 & 0 & -1
\end{vmatrix} = 6,\;
D_3 = \begin{vmatrix}
2 & 1 & 1 \\
1 & -1 & 2 \\
3 & 2 & 0
\end{vmatrix} = 3
\]
故 $x_1 = \frac{-3}{-3} = 1,\; x_2 = \frac{6}{-3} = -2,\; x_3 = \frac{3}{-3} = -1$。

\begin{conclusion}{克莱姆法则的推广}{generalized Cramer's rule}
    对于齐次线性方程组 $\displaystyle \sum_{j=1}^{n} a_{ij}x_j = 0 \; (i=1,2,\ldots,n)$：
    \begin{enumerate}
        \item 若系数行列式 $D \neq 0$，则方程组只有零解（平凡解）$x_1 = x_2 = \cdots = x_n = 0$。
        \item 若方程组有非零解，则必有 $D = 0$。
    \end{enumerate}
    这是判定齐次线性方程组是否有非零解的重要依据。
\end{conclusion}

\subsection{行列式与可逆性的关系}

\begin{conclusion}{方阵可逆的充要条件}{invertibility and determinant}
    设 $A$ 为 $n$ 阶方阵，则以下命题等价：
    \begin{enumerate}
        \item $A$ 可逆（非奇异，满秩）；
        \item $\det(A) \neq 0$；
        \item $A$ 的列（行）向量组线性无关；
        \item $A$ 可以通过初等行变换化为单位矩阵 $I_n$；
        \item 齐次线性方程组 $Ax = 0$ 只有零解；
        \item 对任意 $b$，非齐次线性方程组 $Ax = b$ 有唯一解。
    \end{enumerate}
\end{conclusion}

\begin{formula}{伴随矩阵求逆公式}{adjugate matrix inverse}
    设 $A$ 为 $n$ 阶方阵，$A_{ij}$ 为 $a_{ij}$ 的代数余子式，则
    \[
    A^* = \begin{pmatrix}
    A_{11} & A_{21} & \cdots & A_{n1} \\
    A_{12} & A_{22} & \cdots & A_{n2} \\
    \vdots & \vdots & \ddots & \vdots \\
    A_{1n} & A_{2n} & \cdots & A_{nn}
    \end{pmatrix}
    \]
    称为 $A$ 的\textbf{伴随矩阵}（注意：$A^*$ 中第 $i$ 行第 $j$ 列元素是 $A_{ji}$）。\\
    满足：$A A^* = A^* A = \det(A) \cdot I_n$。\\
    当 $\det(A) \neq 0$ 时，有
    \[
    A^{-1} = \frac{1}{\det(A)} A^*
    \]
\end{formula}

\begin{formula}{二阶矩阵求逆的简便公式}{inverse of 2x2 matrix}
    对于二阶矩阵 $A = \begin{pmatrix} a & b \\ c & d \end{pmatrix}$，若 $\det(A) = ad - bc \neq 0$，则
    \[
    A^{-1} = \frac{1}{ad - bc} \begin{pmatrix} d & -b \\ -c & a \end{pmatrix}
    \]
\end{formula}

\boxed{\textbf{例}} 判断矩阵 $A = \begin{pmatrix} 1 & 2 & 3 \\ 2 & 3 & 1 \\ 3 & 1 & 2 \end{pmatrix}$ 是否可逆，若可逆求其逆矩阵。\\
先求行列式：
\[
\det(A) = \begin{vmatrix}
1 & 2 & 3 \\
2 & 3 & 1 \\
3 & 1 & 2
\end{vmatrix}
= 1(6-1) - 2(4-3) + 3(2-9) = 5 - 2 - 21 = -18 \neq 0
\]
故 $A$ 可逆。计算代数余子式：
\[
\begin{aligned}
A_{11} &= +\begin{vmatrix}3&1\\1&2\end{vmatrix}=5, &
A_{12} &= -\begin{vmatrix}2&1\\3&2\end{vmatrix}=-1, &
A_{13} &= +\begin{vmatrix}2&3\\3&1\end{vmatrix}=-7,\\
A_{21} &= -\begin{vmatrix}2&3\\1&2\end{vmatrix}=-1, &
A_{22} &= +\begin{vmatrix}1&3\\3&2\end{vmatrix}=-7, &
A_{23} &= -\begin{vmatrix}1&2\\3&1\end{vmatrix}=5,\\
A_{31} &= +\begin{vmatrix}2&3\\3&1\end{vmatrix}=-7, &
A_{32} &= -\begin{vmatrix}1&3\\2&1\end{vmatrix}=5, &
A_{33} &= +\begin{vmatrix}1&2\\2&3\end{vmatrix}=-1.
\end{aligned}
\]
故 $A^{-1} = -\frac{1}{18}\begin{pmatrix} 5 & -1 & -7 \\ -1 & -7 & 5 \\ -7 & 5 & -1 \end{pmatrix}$。

\begin{conclusion}{伴随矩阵的性质}{properties of adjugate matrix}
    \begin{enumerate}
        \item $\det(A^*) = [\det(A)]^{n-1}$（当 $n \geq 2$ 时）
        \item $(A^*)^{-1} = \frac{1}{\det(A)}A = (A^{-1})^*$（当 $A$ 可逆时）
        \item $(A^*)^{T} = (A^{T})^*$
        \item $(kA)^* = k^{n-1} A^*$
        \item $(AB)^* = B^* A^*$
        \item $(A^*)^* = [\det(A)]^{n-2} A$（当 $n \geq 3$ 时）
    \end{enumerate}
\end{conclusion}

\subsection{行列式的计算技巧}

\begin{formula}{行列式计算常用方法}{common determinant computation techniques}
    \begin{enumerate}
        \item \textbf{化为三角行列式法}：利用初等行变换（倍加变换为主），将行列式化为上三角形式，
        行列式的值即为主对角线元素的乘积。
        \item \textbf{按行（列）展开法}（降阶法）：选择零元素较多的行或列进行展开，结合性质先制造更多的零。
        \item \textbf{加边法}（升阶法）：通过添加一行一列将原行列式扩大一阶，便于使用性质化简。
        \item \textbf{递推法}：对于具有规律性的 $n$ 阶行列式，建立 $D_n$ 与 $D_{n-1}$（或 $D_{n-2}$）之间的递推关系。
        \item \textbf{数学归纳法}：先计算低阶情形猜测结果，再用归纳法证明。
        \item \textbf{拆项法}：利用行列式的线性性质，将某行（列）拆分为两部分。
        \item \textbf{提取公因子法}：将各行（列）的公因子提出，简化计算。
    \end{enumerate}
\end{formula}

\boxed{\textbf{例}}（加边法）计算 $n$ 阶行列式
\[
D_n = \begin{vmatrix}
1+a_1 & 1 & \cdots & 1 \\
1 & 1+a_2 & \cdots & 1 \\
\vdots & \vdots & \ddots & \vdots \\
1 & 1 & \cdots & 1+a_n
\end{vmatrix}
\]
加边：
\[
D_n = \begin{vmatrix}
1 & 1 & 1 & \cdots & 1 \\
0 & 1+a_1 & 1 & \cdots & 1 \\
0 & 1 & 1+a_2 & \cdots & 1 \\
\vdots & \vdots & \vdots & \ddots & \vdots \\
0 & 1 & 1 & \cdots & 1+a_n
\end{vmatrix}
\xrightarrow{r_i - r_1\;(i\geq 2)}
\begin{vmatrix}
1 & 1 & 1 & \cdots & 1 \\
-1 & a_1 & 0 & \cdots & 0 \\
-1 & 0 & a_2 & \cdots & 0 \\
\vdots & \vdots & \vdots & \ddots & \vdots \\
-1 & 0 & 0 & \cdots & a_n
\end{vmatrix}
\]
再按第一列展开或用分块公式，最终得
\[
D_n = a_1a_2\cdots a_n \left(1 + \sum_{i=1}^{n} \frac{1}{a_i}\right) \quad (a_i \neq 0)
\]

\newpage
\section{习题}

\subsection{基础过关题}

\begin{enumerate}

\item \textbf{填空题} 计算四阶行列式：
\[
D = \begin{vmatrix}
2 & 1 & -1 & 3 \\
4 & 2 & 1 & 5 \\
1 & 0 & 2 & -1 \\
3 & -2 & 1 & 4
\end{vmatrix}
\]
通过初等行变换化为上三角行列式，$D = \underline{\qquad}$。

\ifshowsolutions%
\boxed{\textbf{解}} 利用初等行变换将行列式化为上三角形式。
\[
\begin{aligned}
D &= \begin{vmatrix}
2 & 1 & -1 & 3 \\
4 & 2 & 1 & 5 \\
1 & 0 & 2 & -1 \\
3 & -2 & 1 & 4
\end{vmatrix}
\xrightarrow{r_2 - 2r_1}
\begin{vmatrix}
2 & 1 & -1 & 3 \\
0 & 0 & 3 & -1 \\
1 & 0 & 2 & -1 \\
3 & -2 & 1 & 4
\end{vmatrix} \\
&\xrightarrow{\text{交换 } r_2 \leftrightarrow r_3}
-\begin{vmatrix}
2 & 1 & -1 & 3 \\
1 & 0 & 2 & -1 \\
0 & 0 & 3 & -1 \\
3 & -2 & 1 & 4
\end{vmatrix}
\xrightarrow{r_2 - \frac{1}{2}r_1}
-\begin{vmatrix}
2 & 1 & -1 & 3 \\
0 & -\frac{1}{2} & \frac{5}{2} & -\frac{5}{2} \\
0 & 0 & 3 & -1 \\
3 & -2 & 1 & 4
\end{vmatrix} \\
&\xrightarrow{r_4 - \frac{3}{2}r_1}
-\begin{vmatrix}
2 & 1 & -1 & 3 \\
0 & -\frac{1}{2} & \frac{5}{2} & -\frac{5}{2} \\
0 & 0 & 3 & -1 \\
0 & -\frac{7}{2} & \frac{5}{2} & -\frac{1}{2}
\end{vmatrix}
\xrightarrow{r_4 - 7r_2}
-\begin{vmatrix}
2 & 1 & -1 & 3 \\
0 & -\frac{1}{2} & \frac{5}{2} & -\frac{5}{2} \\
0 & 0 & 3 & -1 \\
0 & 0 & -15 & 17
\end{vmatrix} \\
&\xrightarrow{r_4 + 5r_3}
-\begin{vmatrix}
2 & 1 & -1 & 3 \\
0 & -\frac{1}{2} & \frac{5}{2} & -\frac{5}{2} \\
0 & 0 & 3 & -1 \\
0 & 0 & 0 & 12
\end{vmatrix}
= -\left[2 \cdot \left(-\frac{1}{2}\right) \cdot 3 \cdot 12\right] = -(-36) = 36.
\end{aligned}
\]
\boxed{\textbf{答}} $D = 36$。

\fi%
\item \textbf{选择题} 下列哪种操作会使行列式变为原行列式的\textbf{相反数}？
\begin{enumerate}[label=(\Alph*)]
    \item 交换两行
    \item 将某一行的$k$倍加到另一行上
    \item 将某一行乘以常数 $k$（$k\neq0$）
    \item 将行列式转置
\end{enumerate}

\ifshowsolutions%
\boxed{\textbf{解}} 交换行列式的两行，行列式变为原来的相反数；
倍加变换和转置不改变行列式，将一行乘以 $k$ 则使行列式乘以 $k$。

\boxed{\textbf{答}} 选 \textbf{(A)}。

\fi%
\item \textbf{计算题} 利用按行（列）展开定理，选择零元素最多的行或列，计算：
\[
D = \begin{vmatrix}
3 & 0 & 2 & 1 \\
5 & 2 & 0 & 4 \\
0 & 0 & -1 & 3 \\
2 & 0 & 1 & 0
\end{vmatrix}
\]

\ifshowsolutions%
\boxed{\textbf{解}} 观察行列式，第2列除 $a_{22}=2$ 外其余均为 $0$，选择按第2列展开：
\[
\begin{aligned}
D &= 2 \cdot (-1)^{2+2} \begin{vmatrix}
3 & 2 & 1 \\
0 & -1 & 3 \\
2 & 1 & 0
\end{vmatrix} \quad (\text{第2列其余元素为0，只有一项}) \\
&= 2 \cdot \begin{vmatrix}
3 & 2 & 1 \\
0 & -1 & 3 \\
2 & 1 & 0
\end{vmatrix}
\end{aligned}
\]
将上面的三阶行列式按第1列展开（含一个零）：
\[
\begin{aligned}
\begin{vmatrix}
3 & 2 & 1 \\
0 & -1 & 3 \\
2 & 1 & 0
\end{vmatrix}
&= 3 \cdot (-1)^{1+1} \begin{vmatrix} -1 & 3 \\ 1 & 0 \end{vmatrix}
+ 0 \cdot (\cdots)
+ 2 \cdot (-1)^{3+1} \begin{vmatrix} 2 & 1 \\ -1 & 3 \end{vmatrix} \\
&= 3 \cdot [(-1)\cdot 0 - 3\cdot 1] + 2 \cdot [2\cdot 3 - 1\cdot(-1)] \\
&= 3 \cdot (-3) + 2 \cdot (6 + 1) = -9 + 14 = 5.
\end{aligned}
\]
故 $D = 2 \times 5 = 10$。

\boxed{\textbf{答}} $D = 10$。

\fi%
\item \textbf{计算题} 解关于 $x$ 的方程：
\[
\begin{vmatrix}
x-1 & 2 & 0 \\
2 & x-1 & 2 \\
0 & 2 & x-1
\end{vmatrix} = 0
\]

\ifshowsolutions%
\boxed{\textbf{解}} 计算三阶行列式：
\[
\begin{aligned}
\begin{vmatrix}
x-1 & 2 & 0 \\
2 & x-1 & 2 \\
0 & 2 & x-1
\end{vmatrix}
&= (x-1)\begin{vmatrix} x-1 & 2 \\ 2 & x-1 \end{vmatrix}
- 2\begin{vmatrix} 2 & 2 \\ 0 & x-1 \end{vmatrix} + 0 \\
&= (x-1)[(x-1)^2 - 4] - 2[2(x-1) - 0] \\
&= (x-1)[(x-1)^2 - 4] - 4(x-1) \\
&= (x-1)[(x-1)^2 - 8] = 0.
\end{aligned}
\]
故 $(x-1)[(x-1)^2 - 8] = 0$，即 $x-1=0$ 或 $(x-1)^2=8$，得：
\[
x = 1,\quad x = 1 \pm 2\sqrt{2}.
\]

\boxed{\textbf{答}} $x = 1,\; 1+\sqrt{8},\; 1-\sqrt{8}$。

\fi%
\item \textbf{选择题} 设 $A$ 为 $n$ 阶方阵，$k$ 为常数，则下列等式中正确的是
\begin{enumerate}[label=(\Alph*)]
    \item $|kA| = k|A|$
    \item $|kA| = |k|\cdot|A|$
    \item $|kA| = k^n|A|$
    \item $|kA| = |k|^n\cdot |A|$
\end{enumerate}

\ifshowsolutions%
\boxed{\textbf{解}} 由行列式的公因子提取性质：若将方阵 $A$ 的\textbf{每一行}都乘以 $k$，相当于从每一行提取公因子 $k$，共 $n$ 行，因此 $|kA| = k^n|A|$。\\
选项 (C) 正确。注意 $|kA|$ 中 $k$ 是实数，不需加绝对值。(A) 丢掉了指数 $n$；(B) 多加了绝对值；(D) 错误。常见考研陷阱题。

\boxed{\textbf{答}} 选\textbf{(C)}。

\fi%
\item \textbf{计算题} 利用范德蒙德行列式公式，计算：
\[
V = \begin{vmatrix}
1 & 1 & 1 & 1 \\
1 & 2 & -1 & 3 \\
1 & 4 & 1 & 9 \\
1 & 8 & -1 & 27
\end{vmatrix}
\]

\ifshowsolutions%
\boxed{\textbf{解}} 该行列式第1行为 $(1,1,1,1)$，第2行为 $(x_1,x_2,x_3,x_4)$，第3行为 $(x_1^2,x_2^2,x_3^2,x_4^2)$，第4行为 $(x_1^3,x_2^3,x_3^3,x_4^3)$，这是一个 $4$ 阶范德蒙德行列式，其中
\[
x_1=1,\; x_2=2,\; x_3=-1,\; x_4=3.
\]
由范德蒙德行列式公式：
\[
\begin{aligned}
V &= \prod_{1 \leq i < j \leq 4} (x_j - x_i) \\
&= (2-1)(-1-1)(3-1)(-1-2)(3-2)(3-(-1)) \\
&= 1 \cdot (-2) \cdot 2 \cdot (-3) \cdot 1 \cdot 4 \\
&= 1 \times (-2) \times 2 \times (-3) \times 1 \times 4 = 48.
\end{aligned}
\]

\boxed{\textbf{答}} $V = 48$。

\fi%
\item \textbf{填空题} 设 $A$ 为 $3$ 阶方阵且 $|A|=2$，则
$|2A| = \underline{\qquad}$，
$|A^{-1}| = \underline{\qquad}$，
$|A^*| = \underline{\qquad}$（其中 $A^*$ 为伴随矩阵）。

\ifshowsolutions%
\boxed{\textbf{解}} 已知 $A$ 为 $3$ 阶方阵，$|A|=2$，$n=3$。
\begin{itemize}
    \item $|2A| = 2^3 |A| = 8 \times 2 = 16$（从每一行提取公因子 $2$，共 $3$ 行）。
    \item $|A^{-1}| = |A|^{-1} = \frac{1}{2}$（因为 $AA^{-1}=I$，取行列式得 $|A||A^{-1}|=1$，故 $|A^{-1}|=1/|A|$）。
    \item $|A^*| = |A|^{n-1} = |A|^{2} = 4$（伴随矩阵性质：$|A^*|=|A|^{n-1}$）。
\end{itemize}

\boxed{\textbf{答}} $|2A|=16$，$|A^{-1}|=\frac{1}{2}$，$|A^*|=4$。

\fi%
\end{enumerate}

\subsection{综合提高题}

\begin{enumerate}[resume]

\item \textbf{综合计算题} 设 $A,B$ 均为 $3$ 阶方阵，$|A|=2$，$|B|=-3$，
且满足 $AB + A + B = O$（零矩阵）。求 $|A+B|$。
\boxed{\text{注}} 提示：将原式移项得 $A+B = -AB$，再取行列式。

\ifshowsolutions%
\boxed{\textbf{解}} 由题设直接移项得
\[
A+B=-AB.
\]
因此对三阶方阵取行列式，
\[
|A+B| = |{-AB}| = (-1)^3|AB| = (-1)^3|A||B| = -1 \times 2 \times (-3) = 6.
\]

\boxed{\textbf{答}} $|A+B| = 6$。

\fi%
\item \textbf{证明题} 设 $A$ 为 $n$ 阶方阵（$n \geq 2$），$A^*$ 为 $A$ 的伴随矩阵。
证明：$|A^*| = |A|^{n-1}$。

\ifshowsolutions%
\boxed{\textbf{证明}} 由伴随矩阵的性质，$AA^* = A^*A = |A|I_n$。

若 $|A| \neq 0$（$A$ 可逆）：对 $AA^* = |A|I_n$ 两边取行列式：
\[
|A||A^*| = ||A|I_n| = |A|^n |I_n| = |A|^n.
\]
因为 $|A| \neq 0$，可约去 $|A|$，得 $|A^*| = |A|^{n-1}$。

若 $|A| = 0$：需证 $|A^*| = 0$。用反证法。假设 $|A^*| \neq 0$，则 $A^*$ 可逆。由 $AA^* = |A|I_n = 0 \cdot I_n = O$，右乘 $(A^*)^{-1}$ 得 $A = O$。但若 $A = O$，则 $A^* = O$（零矩阵的伴随矩阵也是零矩阵），与 $|A^*| \neq 0$ 矛盾。故 $|A^*| = 0 = 0^{n-1} = |A|^{n-1}$。

当 $n-1=0$ 即 $n=1$ 时需单独约定，但题设 $n \geq 2$，公式成立。

\boxed{\textbf{证毕}} $|A^*| = |A|^{n-1}$。

\fi%
\item \textbf{计算题}（三对角行列式）计算 $n$ 阶行列式：
\[
D_n = \begin{vmatrix}
2 & 1 & 0 & \cdots & 0 & 0 \\
1 & 2 & 1 & \cdots & 0 & 0 \\
0 & 1 & 2 & \cdots & 0 & 0 \\
\vdots & \vdots & \vdots & \ddots & \vdots & \vdots \\
0 & 0 & 0 & \cdots & 2 & 1 \\
0 & 0 & 0 & \cdots & 1 & 2
\end{vmatrix}
\]
\boxed{\text{注}} 提示：按第一行展开，建立递推关系 $D_n = 2D_{n-1} - D_{n-2}$。

\ifshowsolutions%
\boxed{\textbf{解}} 记 $n$ 阶三对角行列式为 $D_n$。当 $n=1$ 时，$D_1 = |2| = 2$。当 $n=2$ 时，
\[
D_2 = \begin{vmatrix} 2 & 1 \\ 1 & 2 \end{vmatrix} = 3.
\]

对 $n\geq3$，按第一行展开。第一项的余子式就是 $D_{n-1}$；
第二项的余子式记为 $M_{12}$，其第一列只有左上角元素 $1$ 非零，
再按该列展开即得 $M_{12}=D_{n-2}$。考虑代数余子式的符号，故
\[
D_n=2D_{n-1}-D_{n-2}.
\]

因此得到递推公式 \(D_n=2D_{n-1}-D_{n-2}\ (n\geq3)\)。

这是二阶线性递推，特征方程 $r^2 - 2r + 1 = 0$，得 $r=1$（二重根），通解为：
\[
D_n = (c_1 + c_2 n) \cdot 1^n = c_1 + c_2 n.
\]
代入初值 $D_1 = 2$，$D_2 = 3$：
\[
\begin{cases}
c_1 + c_2 = 2, \\
c_1 + 2c_2 = 3,
\end{cases}
\]
解得 $c_2 = 1$，$c_1 = 1$。故 $D_n = n+1$。

\boxed{\textbf{答}} $D_n = n+1$。

\fi%
\item \textbf{证明题} 设 $A$ 为 $n$ 阶反对称矩阵（即 $A^T = -A$），且 $n$ 为奇数。
证明：$|A| = 0$。

\ifshowsolutions%
\boxed{\textbf{证明}} 设 $A$ 为 $n$ 阶反对称矩阵，即 $A^T = -A$，且 $n$ 为奇数。

由行列式的转置不变性：$|A^T| = |A|$。\\
又由 $A^T = -A$，得 $|A^T| = |{-A}|$。

由于 $n$ 阶方阵乘以 $-1$，等于从每一行提取因子 $-1$，共 $n$ 行：
\[
|{-A}| = (-1)^n |A|.
\]

结合以上：
\[
|A| = |A^T| = |{-A}| = (-1)^n |A|.
\]

因为 $n$ 为奇数，$(-1)^n = -1$，因此：
\[
|A| = -|A| \implies 2|A| = 0 \implies |A| = 0.
\]

\boxed{\textbf{证毕}} 奇数阶反对称矩阵的行列式为零。

\fi%
\item \textbf{计算题} 用克莱姆法则求解方程组，并用高斯消元法验证结果：
\[
\begin{cases}
2x_1 + x_2 - x_3 = 5 \\
x_1 - 3x_2 + 2x_3 = -1 \\
3x_1 + 2x_2 + x_3 = 8
\end{cases}
\]

\ifshowsolutions%
\boxed{\textbf{解}} 
\textbf{法一：克莱姆法则。}

系数行列式：
\[
D = \begin{vmatrix}
2 & 1 & -1 \\
1 & -3 & 2 \\
3 & 2 & 1
\end{vmatrix}
\]
按第1行展开：
\[
\begin{aligned}
D &= 2\begin{vmatrix} -3 & 2 \\ 2 & 1 \end{vmatrix}
- 1\begin{vmatrix} 1 & 2 \\ 3 & 1 \end{vmatrix}
+ (-1)\begin{vmatrix} 1 & -3 \\ 3 & 2 \end{vmatrix} \\
&= 2[(-3)\cdot 1 - 2\cdot 2] - 1[1\cdot 1 - 2\cdot 3] - [1\cdot 2 - (-3)\cdot 3] \\
&= 2(-3 - 4) - (1 - 6) - (2 + 9) \\
&= 2(-7) - (-5) - 11 = -14 + 5 - 11 = -20.
\end{aligned}
\]

计算 $D_1$（将第1列换为常数项）：
\[
D_1 = \begin{vmatrix}
5 & 1 & -1 \\
-1 & -3 & 2 \\
8 & 2 & 1
\end{vmatrix}
\]
按第1行展开：
\[
\begin{aligned}
D_1 &= 5\begin{vmatrix} -3 & 2 \\ 2 & 1 \end{vmatrix}
- 1\begin{vmatrix} -1 & 2 \\ 8 & 1 \end{vmatrix}
+ (-1)\begin{vmatrix} -1 & -3 \\ 8 & 2 \end{vmatrix} \\
&= 5(-7) - 1(-1 - 16) - [(-1)\cdot 2 - (-3)\cdot 8] \\
&= -35 - (-17) - (-2 + 24) = -35 + 17 - 22 = -40.
\end{aligned}
\]

计算 $D_2$（将第2列换为常数项）：
\[
D_2 = \begin{vmatrix}
2 & 5 & -1 \\
1 & -1 & 2 \\
3 & 8 & 1
\end{vmatrix}
\]
按第1行展开：
\[
\begin{aligned}
D_2 &= 2\begin{vmatrix} -1 & 2 \\ 8 & 1 \end{vmatrix}
- 5\begin{vmatrix} 1 & 2 \\ 3 & 1 \end{vmatrix}
+ (-1)\begin{vmatrix} 1 & -1 \\ 3 & 8 \end{vmatrix} \\
&= 2(-1 - 16) - 5(1 - 6) - (8 + 3) \\
&= 2(-17) - 5(-5) - 11 = -34 + 25 - 11 = -20.
\end{aligned}
\]

计算 $D_3$（将第3列换为常数项）：
\[
D_3 = \begin{vmatrix}
2 & 1 & 5 \\
1 & -3 & -1 \\
3 & 2 & 8
\end{vmatrix}
\]
按第1行展开：
\[
\begin{aligned}
D_3 &= 2\begin{vmatrix} -3 & -1 \\ 2 & 8 \end{vmatrix}
- 1\begin{vmatrix} 1 & -1 \\ 3 & 8 \end{vmatrix}
+ 5\begin{vmatrix} 1 & -3 \\ 3 & 2 \end{vmatrix} \\
&= 2(-24 + 2) - 1(8 + 3) + 5(2 + 9) \\
&= 2(-22) - 11 + 5\cdot 11 = -44 - 11 + 55 = 0.
\end{aligned}
\]

由克莱姆法则：
\[
x_1 = \frac{D_1}{D} = \frac{-40}{-20} = 2,\quad
x_2 = \frac{D_2}{D} = \frac{-20}{-20} = 1,\quad
x_3 = \frac{D_3}{D} = \frac{0}{-20} = 0.
\]

\textbf{法二：高斯消元法验证。}

增广矩阵：
\[
\begin{pmatrix}
2 & 1 & -1 & 5 \\
1 & -3 & 2 & -1 \\
3 & 2 & 1 & 8
\end{pmatrix}
\xrightarrow{r_2 - \frac{1}{2}r_1}
\begin{pmatrix}
2 & 1 & -1 & 5 \\
0 & -\frac{7}{2} & \frac{5}{2} & -\frac{7}{2} \\
3 & 2 & 1 & 8
\end{pmatrix}
\xrightarrow{r_3 - \frac{3}{2}r_1}
\begin{pmatrix}
2 & 1 & -1 & 5 \\
0 & -\frac{7}{2} & \frac{5}{2} & -\frac{7}{2} \\
0 & \frac{1}{2} & \frac{5}{2} & \frac{1}{2}
\end{pmatrix}
\]
\[
\xrightarrow{r_3 + \frac{1}{7}r_2}
\begin{pmatrix}
2 & 1 & -1 & 5 \\
0 & -\frac{7}{2} & \frac{5}{2} & -\frac{7}{2} \\
0 & 0 & \frac{20}{7} & 0
\end{pmatrix}
\]
回代：$\frac{20}{7}x_3 = 0 \implies x_3 = 0$；
$-\frac{7}{2}x_2 + \frac{5}{2}x_3 = -\frac{7}{2} \implies -\frac{7}{2}x_2 = -\frac{7}{2} \implies x_2 = 1$；
$2x_1 + x_2 - x_3 = 5 \implies 2x_1 + 1 = 5 \implies x_1 = 2$。
结果一致。

\boxed{\textbf{答}} $x_1 = 2,\; x_2 = 1,\; x_3 = 0$。

\fi%
\end{enumerate}

\subsection{真题风格与综合训练}

\begin{enumerate}[resume]

\item \textbf{填空题} $\boxed{\textbf{真题风格}}$ 设 $\alpha_1,\alpha_2,\alpha_3$ 为三维列向量，记矩阵
$A = (\alpha_1,\alpha_2,\alpha_3)$，$B = (\alpha_1+2\alpha_2,\;\alpha_2+2\alpha_3,\;\alpha_3+2\alpha_1)$。
若 $|A| = 3$，则 $|B| = \underline{\qquad}$。

\ifshowsolutions%
\boxed{\textbf{解}} 由已知：
\[
B = (\alpha_1+2\alpha_2,\;\alpha_2+2\alpha_3,\;\alpha_3+2\alpha_1).
\]

矩阵 $B$ 可由 $A$ 经列变换得到：
\[
B = (\alpha_1,\alpha_2,\alpha_3) \begin{pmatrix}
1 & 0 & 2 \\
2 & 1 & 0 \\
0 & 2 & 1
\end{pmatrix}
= A \cdot C,
\]
其中
\[
C = \begin{pmatrix}
1 & 0 & 2 \\
2 & 1 & 0 \\
0 & 2 & 1
\end{pmatrix}.
\]

取行列式：
\[
|B| = |A| \cdot |C|.
\]

计算 $|C|$：
\[
\begin{aligned}
|C| &= \begin{vmatrix}
1 & 0 & 2 \\
2 & 1 & 0 \\
0 & 2 & 1
\end{vmatrix}
= 1\begin{vmatrix} 1 & 0 \\ 2 & 1 \end{vmatrix} - 0 + 2\begin{vmatrix} 2 & 1 \\ 0 & 2 \end{vmatrix} \\
&= 1 \cdot 1 + 2 \cdot 4 = 9.
\end{aligned}
\]

故 $|B| = 3 \times 9 = 27$。

\boxed{\textbf{答}} $|B| = 27$。

\fi%
\item \textbf{计算题} $\boxed{\textbf{真题风格}}$ 计算 $n$ 阶行列式：
\[
D_n = \begin{vmatrix}
a & b & b & \cdots & b \\
b & a & b & \cdots & b \\
b & b & a & \cdots & b \\
\vdots & \vdots & \vdots & \ddots & \vdots \\
b & b & b & \cdots & a
\end{vmatrix}_{n\times n}
\]
即主对角线元素均为 $a$，其余元素均为 $b$。

\ifshowsolutions%
\boxed{\textbf{解}} 将第2至第 $n$ 列都加到第1列：
\[
D_n = \begin{vmatrix}
a+(n-1)b & b & b & \cdots & b \\
a+(n-1)b & a & b & \cdots & b \\
a+(n-1)b & b & a & \cdots & b \\
\vdots & \vdots & \vdots & \ddots & \vdots \\
a+(n-1)b & b & b & \cdots & a
\end{vmatrix}
\]

提取第1列公因子 $[a+(n-1)b]$：
\[
D_n = [a+(n-1)b] \begin{vmatrix}
1 & b & b & \cdots & b \\
1 & a & b & \cdots & b \\
1 & b & a & \cdots & b \\
\vdots & \vdots & \vdots & \ddots & \vdots \\
1 & b & b & \cdots & a
\end{vmatrix}
\]

将第1行的 $(-b)$ 倍加到其余各行（$i \geq 2$，注意第 $i$ 行的第 $j$ 列，$j \geq 2$）：
\[
D_n = [a+(n-1)b] \begin{vmatrix}
1 & b & b & \cdots & b \\
0 & a-b & 0 & \cdots & 0 \\
0 & 0 & a-b & \cdots & 0 \\
\vdots & \vdots & \vdots & \ddots & \vdots \\
0 & 0 & 0 & \cdots & a-b
\end{vmatrix}
\]

这是上三角行列式（实际是具第一行非对角元的下三角），其值等于主对角线元素之积：
\[
D_n = [a+(n-1)b] \cdot 1 \cdot (a-b)^{n-1} = [a + (n-1)b](a-b)^{n-1}.
\]

\boxed{\textbf{答}} $D_n = [a + (n-1)b](a-b)^{n-1}$。

\fi%
\item \textbf{证明题} $\boxed{\textbf{真题风格}}$ 设 $A$ 为 $n$ 阶方阵（$n \geq 2$），$A^*$ 为 $A$ 的伴随矩阵。
证明：
\[
r(A^*) = \begin{cases}
n, & r(A) = n, \\
1, & r(A) = n-1, \\
0, & r(A) \leq n-2.
\end{cases}
\]
其中 $r(\cdot)$ 表示矩阵的秩。

\ifshowsolutions%
\boxed{\textbf{证明}} 分三种情况讨论。

\textbf{情况一：$r(A) = n$。} 此时 $A$ 满秩，$|A| \neq 0$。由 $AA^* = |A|I_n$，两边取行列式：
\[
|A||A^*| = |A|^n \implies |A^*| = |A|^{n-1} \neq 0.
\]
$|A^*| \neq 0$ 说明 $A^*$ 可逆，故 $r(A^*) = n$。

\textbf{情况二：$r(A) = n-1$。} 此时 $|A| = 0$，故 $AA^* = O$。

由 $AA^* = O$，可知 $A^*$ 的每一列都是齐次线性方程组 $Ax = 0$ 的解。因为 $r(A)=n-1$，$Ax=0$ 的解空间维数为 $1$。故 $A^*$ 的列向量都在这条一维子空间中，因此 $r(A^*) \leq 1$。

下面证明 $r(A^*) \neq 0$（即 $A^* \neq O$）。由于 $r(A)=n-1$，至少存在一个 $n-1$ 阶子式非零（因为秩为 $n-1$ 意味着存在非零的 $n-1$ 阶余子式），从而至少有一个代数余子式 $A_{ij} \neq 0$。而 $A^*$ 的元素就是这些代数余子式（转置后），故 $A^* \neq O$，$r(A^*) \geq 1$。

结合 $r(A^*) \leq 1$，得 $r(A^*) = 1$。

\textbf{情况三：$r(A) \leq n-2$。} 此时 $A$ 的所有 $n-1$ 阶子式均为 $0$，故所有代数余子式 $A_{ij} = 0$，从而 $A^* = O$，$r(A^*) = 0$。

综上所述：
\[
r(A^*) = \begin{cases}
n, & r(A)=n, \\
1, & r(A)=n-1, \\
0, & r(A) \leq n-2.
\end{cases}
\]

\boxed{\textbf{证毕}}

\fi%
\end{enumerate}
